\begin{textsample}{POS Dim 4 – pb – Score 16.00 – pb/pb\_inf\_4.txt}  \label{ex:f4_pos_161}
3.3.2 As crianças da Paraíba : \textbf{populações}, identidades, currículo e práticas pedagógicas Para pensar na criança real com a qual os/as professores/as de Educação Infantil da Paraíba vão conviver, é importante conhecer os pertencimentos sociais e culturais desta criança.
A \textbf{população} paraibana é composta, segundo o último censo, por 3.766.528 habitantes, distribuídos em uma área territorial de 56.468,435 km².
A maior parte dessa \textbf{população} está situada nas cidades de João Pessoa e Campina Grande, correspondendo a 40 \% da \textbf{população} do estado.
Do ponto de vista da área, 60 \% dessa \textbf{população} está concentrada em áreas urbanas e 40 \% situam -se em áreas rurais. ( IBGE, 2010 ).
O povo paraibano é constituído pela miscigenação de \textbf{brancos} europeus, \textbf{índios} nativos e negros africanos.
Apesar desse grande entrelaçamento \textbf{racial}, evidencia -se a predominância de alguns grupos étnicos em determinadas microrregiões, a exemplo dos \textbf{indígenas} na Baía da Traição.
Os mesmos realizam suas práticas artesanais, religiosas e alimentares, vivenciando princípios de seu povo.
Há várias comunidades quilombolas situadas do \textbf{litoral} ao \textbf{sertão}.
Estas também vivenciam suas atividades culturais, desde as danças de cunho festivo aos rituais sagrados, tendo a alimentação em consonância com as oferendas aos orixás, em que a arte dos adereços está associada aos elementos ritualísticos, sendo o batuque uma mola mestre da cultura desse povo. ( IBGE, 2010 ).
Há também uma parcela da \textbf{população} de descendência europeia, que vive principalmente nos grandes centros urbanos do estado, protagonizando em seus costumes o cristianismo e as ações vistas como eurocêntricas, evidentes nas crenças aos santos e nos ritos cristãos, distinguindo o sagrado do profano como divisões das ações socioculturais.
Compõem ainda a \textbf{população} paraibana os mulatos, que estão predominantes no \textbf{litoral} \textbf{sul} e no \textbf{agreste}, e os caboclos, que estão mais presentes em todo o interior e no \textbf{litoral} \textbf{norte}. ( IBGE, 2010 ).
Entre os povos que habitam os \textbf{territórios} paraibanos, há também os \textbf{ciganos}, que têm a quiromancia e a leitura de mão como elementos relacionados às crenças dos povos antigos e que aqui se estabeleceram em várias regiões do nosso estado, sobretudo em Sousa.
Também há a presença de comunidades de pescadores, que estão ligados diretamente às regiões alagadiças, sejam \textbf{mar}, \textbf{rios} ou ribeirinhas, que têm nas suas práticas semelhanças, mas que, no convívio mais próximo com sua cultura, percebemos diferenças entre elas.
Mapa dos povos paraibanos Mapa criado de acordo com as fontes citadas no box deste texto.
Para saber mais sobre os povos paraibanos, acesse : - http://www.periodicos.ufpb.br/index.php/ci/article/viewfile/27579/14917 Comunidades quilombolas na Paraíba.
Maria Ester Fortes ( Instituto Nacional de Colonização e Reforma Agrária – Incra ).
Cadernos Imbondeiro.
João Pessoa, v. 4, n. 1, out. 2015. - http://www.ccen.ufpb.br/ppgg/contents/documentos/teses/comunidades-tradicionais-de-pesca-artesanal-marinha-na-paraiba\_realidade-e-desafios.pdf - http://www.ch.ufcg.edu.br/arius/01\_revistas/v19n1/00\_arius\_v19\_n1\_2013\_edicao\_completa.pdf - http://salutte.com.br/aventura/wp-content/uploads/2011/11/mapa-tur\%c3\%adstico-da-para\%c3\%adba.jpg - http://paraiba.pb.gov.br/iphaep-debate-historia-e-resistencia-do-povo-cigano-em-escola-da-capital Mapas de localização do estado da Paraíba no Brasil e dos municípios visitados em relação ao \textbf{litoral} paraibano.
Fonte : IBGE ( 2016 ) ; Ismael Araújo.
Página 52.
Esse mapeamento dos povos paraibanos evidencia traços da sua realidade, contribuindo para que possamos olhar para as diversas infâncias que fazem parte desses povos.
É importante reconhecer que ainda faltam fontes de informações mais sistematizadas para que se possa trazer um cenário mais fiel de todo o povo paraibano, mas esse mapa já dá mostras da rica diversidade de povos que constituem o estado da Paraíba, não podendo eles ser desconsiderados quando se pensa na oferta de Educação Infantil no estado e na relação que esta deve ter quanto ao reconhecimento do pertencimento cultural das crianças.
Assim, é importante reconhecer a distribuição da \textbf{população} paraibana no espaço geográfico, uma vez que a criança se desenvolve e se constitui em contextos diversos, sendo necessárias práticas pedagógicas que reconheçam seus pertencimentos culturais, conforme determinam as Diretrizes Curriculares Nacionais para Educação Infantil ( BRASIL, 2009 ), em seu Artigo 8º.
Em seu livro Reflexões : a criança, o brinquedo e a educação, Walter Benjamin ( 2002 ) ressalta que, para a criança, a brincadeira é libertação.
Para ele, “ brincar significa sempre libertação.
Rodeada por um mundo de gigantes, as crianças criam para si, brincando, um mundo próprio. ” ( BENJAMIN, 2002, p. 64 ).
A organização curricular da Educação Infantil pode se estruturar em eixos, centros, campos ou módulos de experiências que devem se articular em torno dos princípios, condições e objetivos propostos nesta diretriz. ( Parecer CNE, 2009, p. 16 ). > Art. 8º.
A proposta pedagógica das instituições de Educação Infantil deve ter como objetivo garantir à criança acesso a processos de apropriação, renovação e articulação de conhecimentos e aprendizagens de diferentes linguagens, assim como o direito à proteção, à saúde, à liberdade, à confiança, ao respeito, à dignidade, à brincadeira, à convivência e à interação com outras crianças. ( BRASIL, 2009 ).
Esse reconhecimento do pertencimento cultural das crianças e de suas famílias tem efeitos no currículo a ser construído e organizado pela instituição educacional, seus professores, crianças e suas famílias, no sentido de organizá -lo de forma a garantir ações que proporcionem experiências e interações voltadas para o desenvolvimento integral da criança, concebendo -a em todo o processo como sujeito de direitos, capaz de se expressar desde cedo e que, nas interações com os adultos e seus pares, não apenas recebe influência, mas influencia, constrói e se reconstrói na relação com o outro.
Sendo assim, ao pensarmos na organização do currículo, é preciso levar em consideração tudo aquilo que envolve o desenvolvimento da criança, como sujeito social, histórico e cultural.
A Base Nacional Comum Curricular para a Educação Infantil ( BRASIL, 2017 ) elegeu campos de experiência como o modo de organizar o currículo nesta etapa da Educação.
Os campos de experiência — o eu, o outro e o nós ; traços, sons, cores e formas ; escuta, fala, pensamento e imaginação ; corpo, gestos e movimento ; e espaços, tempos, quantidades, relações e transformações — devem ser pensados de maneira articulada e ser explorados a partir dos interesses das crianças.
Eles estão diretamente ligados aos dois eixos curriculares que fundamentam as práticas pedagógicas na Educação Infantil : as interações e brincadeiras.
Esses eixos já estão normatizados nas Diretrizes Curriculares Nacionais para a Educação Infantil ( BRASIL, 2009 ).
Tomando a Educação Infantil como um dos momentos em que a criança começa a ampliar suas relações sociais, as interações e brincadeiras articulam sujeitos e ações e permitem à criança a construção e reconstrução de realidades nas ações que realiza.
Os/As professores/as que atuam nas instituições de Educação Infantil deverão organizar práticas que garantam experiências significativas, levando em consideração a organização dos tempos e espaços, a construção dos materiais, a definição das atividades pedagógicas, com práticas que garantam o desenvolvimento pleno das crianças.
Pensando nas crianças com as quais convive ou que fazem parte da sua história, você consegue identificar situações em que essas crianças se expressam livremente?
Também, em algum momento, já se deparou com um relato de uma mãe no qual ela se surpreende com a criança realizando ações próprias ao adulto?
Ou observou ações da criança parecidas com a observação desta mãe?
Quem nunca ficou encantado ao observar o bebê que brinca com as mãozinhas ou com o barulho do seu balbucio?
Essas situações, carregadas de experiências e aprendizados, expressam momentos nos quais é possível perceber a criança em desenvolvimento, seja de maneira livre ou intencional, seja ainda realizando ações próprias à cultura infantil ou relacionadas à cultura do adulto.
A criança brinca de diversas formas, com o corpo, o movimento, por meio de jogos de regras, jogos simbólicos e faz de conta.
Tendo clareza da importância do brincar, os/as professores/as da Educação Infantil devem tê -lo como eixo fundamental das práticas realizadas junto às crianças.
Outro aspecto a ser considerado no processo de organização curricular diz respeito aos princípios que devem norteá -la nas creches e pré-escolas.
Estes princípios, reconhecidos como sendo de caráter ético, estético e político ( BRASIL, 2009 ), são a base para se pensar os direitos de aprendizagem, como veremos adiante.
Oliveira, refletindo sobre os princípios norteadores da Educação Infantil na relação com os seus direitos, ressalta que as práticas cotidianas devem garantir às crianças o direito de viver a infância protegida, livre, com dignidade e em convivência com outras crianças. ( OLIVEIRA, 2010 ).
Tais princípios fundamentam -se na formação integral do ser humano que, na construção de sua identidade pessoal e social, deve propiciar experiências fundamentadas em princípios éticos, estéticos e políticos.
Os direitos de aprendizagem, definidos na BNCC ( 2017 ), estão diretamente vinculados a esses princípios.
Os direitos de conhecer -se e de conviver remetem aos princípios éticos, pois estão relacionados às ações e às relações estabelecidas com e entre criança/criança, criança/adulto e adulto/criança.
Os princípios políticos estão ligados aos direitos de expressar e de participar.
Já os princípios estéticos manifestam -se nos direitos de brincar e explorar. > “ Brincar de forma livre e prazerosa permite que a criança seja conduzida a uma esfera imaginária, um mundo de faz de conta consciente, porém capaz de reproduzir as relações que observa em seu cotidiano, vivenciando simbolicamente diferentes papéis, exercitando sua capacidade de generalizar e abstrair ” ( MELO \& VALLE, 2005, p. 45 ).
Princípios : - Éticos : da autonomia, da responsabilidade, da solidariedade e do respeito ao bem comum, ao meio ambiente e às diferentes culturas, identidades e singularidades. - Políticos : dos direitos de cidadania, do exercício da criticidade e do respeito à ordem democrática. - Estéticos : da sensibilidade, da criatividade, da ludicidade e da liberdade de expressão nas diferentes manifestações artísticas e culturais. ( BRASIL, 2009 ).
Nas últimas décadas, vem se consolidando, na Educação Infantil, a concepção que vincula educar e cuidar, entendendo o cuidado como algo indissociável do processo educativo.
Nesse contexto, as creches e pré-escolas, ao acolher as vivências e os conhecimentos construídos pelas crianças no ambiente da família e no contexto de sua comunidade, e articulá -los em suas propostas pedagógicas [... ]. ( BRASIL, 2017, p. 34 ).
Em seu primeiro dia na creche, uma criança, ao chegar, recusa -se a entrar e começa a chorar.
A professora, ao recebê -la, coloca -a no colo, começa a acalentá -la e a leva para conhecer o espaço, ao tempo que vai apresentando ela aos coleguinhas.
Aos poucos, a criança vai parando de chorar.
No primeiro dia, a criança pergunta várias vezes se a mamãe vem.
Depois de alguns dias, a criança, já adaptada, não chora mais e diz alegremente para a professora que, depois de determinado momento da rotina, é a hora da chegada da mamãe para pegá -la. ( Relato de experiência em uma creche do município de João Pessoa-PB. ) > “ Cuidar e educar são ações intrínsecas e de responsabilidade da família, dos professores e dos médicos.
Todos têm de saber que só se cuida educando e só se educa cuidando ” ( DIDONET, 2003 ).
Enquanto professores/as em interação cotidiana com as crianças paraibanas no contexto das instituições de Educação Infantil, devemos compreender que nossas ações estão carregadas de intencionalidade, sendo necessária uma avaliação permanente dos processos vivenciados em nossas instituições, tendo como fio condutor nosso compromisso com os princípios que alicerçam práticas junto à criança, tomada como sujeito de direitos, e que, portanto, têm efeitos em sua formação.
Para que a Educação Infantil cumpra seu papel na direção de contribuir para a promoção do desenvolvimento integral da criança, dois aspectos devem ser tomados de forma integrada e indissociável : o cuidar e o educar.
A história nos mostra que o atendimento à criança pequena em creche foi criado com um caráter assistencialista, sendo esta muitas vezes uma instituição emergencial e de guarda para crianças pobres.
Nesse contexto, servia apenas para crianças carentes, filhas de mães trabalhadoras, cumprindo um papel predominantemente vinculado ao cuidado, embora saibamos que o ato educativo também está presente no ato de cuidar.
A Constituição Federal ( BRASIL, 1988 ) e a Lei de Diretrizes e Bases da Educação Nacional/LDB ( BRASIL, 2013 ) contribuíram para pensarmos a Educação Infantil em outro patamar.
Além de ser ofertada em creches e pré-escolas, passou a ser considerada a primeira etapa da Educação Básica.
Esse avanço na Lei produz condições para mudanças nas concepções e práticas de todos que atuam em creches e pré-escolas.
As Diretrizes Curriculares Nacionais para a Educação Infantil ( BRASIL, 2009 ) e a Base Nacional Comum Curricular ( BRASIL, 2017 ) vêm reforçar a importância de um olhar para a infância respeitando suas especificidades e garantindo seus direitos.
As ações de cuidar e educar devem estar integradas, sendo função e papel das instituições de Educação Infantil favorecer o desenvolvimento da criança como um ser completo, não sendo possível educar sem cuidar e vice-versa.
Nas rotinas diárias, desde o momento que a criança chega às creches e pré-escolas, o cuidar e o educar vão acontecendo simultaneamente, favorecendo à criança o seu desenvolvimento integral, no que tange aos aspectos cognitivos, afetivos, emocionais, sociais, físicos, estéticos e éticos.
Esse relato mostra as ações de cuidar e educar ocorrendo simultaneamente.
A acolhida da professora, sua sensibilidade e empatia ocorreram ao mesmo tempo em que ajudou a criança a se localizar no novo contexto, tanto quanto ao espaço como quanto às pessoas.
A adaptação da criança demonstra que ela compreendeu o seu lugar ali, tornando -se íntima ao ambiente e avançando em seu desenvolvimento social.
Essa percepção da criança demonstra uma visão integral da mesma.
Atuar junto à criança de zero a cinco anos de idade apresenta diversos desafios.
Muitas vezes, nos perguntamos como professor/a de creche sobre o que fazer para que as crianças recebam o “ nosso ” melhor.
Na direção que estamos pensando a criança aqui, cabe, sobretudo, perguntar : O que é o melhor também para as crianças ou para cada criança?
Esse é um exercício que não anula o professor, mas o conduz a pensar na alteridade que cada criança representa e na necessidade de ouvir, ver, compreender, enfim, ser sensível a essa alteridade.
Ainda, para responder a questões como essas que tanto nos inquietam, é necessário pensar na educação de forma ampla, entendendo que a criança é um ser integral e que as instituições de Educação Infantil devem favorecer ações que promovam o seu desenvolvimento numa perspectiva que incorpore seus aspectos físicos, emocionais, afetivos, cognitivos, linguísticos e sociais.
Sendo assim, requer posturas e conhecimentos específicos que compreendam a singularidade de cada criança, pois ela está se construindo como sujeito e descobrindo o mundo, fazendo escolhas referentes ao jeito de ser e estar com o outro, construindo sua identidade pessoal e coletiva.
Dessa forma, é função da Educação Infantil oportunizar às crianças vivências que favoreçam o seu desenvolvimento integral em todo o processo de desenvolvimento.
As instituições que se propõem a atuar com a Educação Infantil não podem ser condicionadas a pensar apenas nas linguagens da fala ou da escrita, devendo dar importância a outras linguagens e ações, como o movimento, a brincadeira, o desenho, a dramatização, a dança, a música, o gesto, entre outros.
Ou seja, devem se constituir em espaços ricos e potentes nos quais as crianças possam ter seus direitos respeitados.

% matched lemmas: agreste, branco, cigano, indígena, litoral, mar, norte, população, racial, rio, sertão, sul, território, índio
\end{textsample}
